% ID: SEL-D0024-OA-E02
% ESTADO: cotejada. Segunda pasada visual completa y borrador compilado revisados.
% Fuente: Selectividad/matematicas_aplicadas/reserva_b_EXAMEN 2-ANDALUCÍA.pdf
% SHA-256: 98b7723232acd3143de28397f7641923ed48292db7c35883ebd2b56a72b30b79
% Página: 1; opción A; ejercicio original 2.

Una empresa quiere invertir en productos financieros un mínimo de un millón de euros y un máximo de seis millones de euros. La rentabilidad que obtiene viene dada en función de la cantidad invertida, $x$, por la siguiente expresión:
\[
R(x)=\begin{cases}
x-2 & \text{si }1\leq x<2\\[4pt]
-x^2+10x-16 & \text{si }2\leq x\leq 6
\end{cases}
\]
donde tanto $x$, como $R(x)$, están expresadas en millones de euros.
\begin{enumerate}[label=\alph*)]
\item \textbf{(0.75 puntos)} Estudie la continuidad de la función $R$.
\item \textbf{(0.75 puntos)} Esboce la gráfica de la función.
\item \textbf{(1 punto)} \textquestiondown{}Qué cantidad debe invertir para obtener la máxima rentabilidad y a cuánto asciende ésta? \textquestiondown{}Para qué valores de $x$ la rentabilidad es positiva?
\end{enumerate}
